# Related Work — Parts 2 and 3

## 2. Graph Self-Supervised and Multi-View Learning

Self-supervised learning (SSL) is an appealing alternative to imputation on
attribute-missing graphs. Instead of filling in absent inputs, an SSL method learns
representations from the unlabeled structure and is judged by how well those
representations support downstream tasks. Our work belongs to this line, so we review it in
some detail. We keep apart two strands that are often merged. Technique neighbors are
general graph SSL methods developed on fully observed graphs, and we borrow our objectives
and views from them. Problem neighbors are SSL methods designed for attribute-missing
graphs, and these are the ones we must separate ourselves from most carefully.

### 2.1 Contrastive and multi-view self-supervision

The dominant approach learns node representations by making different views of a graph
agree, usually framed as maximizing a lower bound on their mutual information. Methods in
this family differ in the granularity of the agreement and in how the views are produced.
Deep Graph Infomax [16] maximizes the mutual information between patch-level node
representations and a graph-level summary, avoids random-walk objectives, and works in both
transductive and inductive settings. GMI [17] measures mutual information between a node's
hidden representation and its neighborhood over both features and topology, which ties the
signal to local structure. MVGRL [15] contrasts an adjacency view against a graph-diffusion
view. Its authors report that adding more than two views or contrasting multi-scale
encodings does not help, and that the strongest results come from contrasting first-order
neighbors with the diffusion view.

A large group of methods build views by perturbing the input at random. GRACE [18] corrupts
structure and attributes to form two views and contrasts them node by node. GCA [19] makes
the corruption adaptive, perturbing less important edges and features more aggressively
according to centrality and attribute-importance scores. GraphCL [21] catalogs four families
of graph augmentation and studies how they combine, and later work automates the choice
because it otherwise has to be retuned for each dataset. Handcrafted augmentation can change
what a graph means, so a parallel line reduces or removes it. AFGRL [25] forms a second view
without augmentation by finding nodes that share local structure and global semantics with
the target. MA-GCL [20] perturbs the view encoders' architectures, through asymmetric,
random, and shuffling tricks, rather than the graph itself. A final group drops negative
pairs and decorrelates the embedding dimensions in the spirit of Barlow Twins [24], whose
graph counterpart is DCLN [26].

Taken together, these methods show that the objective and the view both matter, and that
generic augmentations are a brittle design axis. All of them, however, are built and tested
on fully observed graphs. Their views come from augmentations that are not designed for
absent features, and none of them asks which objective still carries information once whole
attribute vectors are gone. That is the question we take up.

### 2.2 Masked and reconstruction-based self-supervision

A second family of pretext tasks corrupts part of the input and trains the model to restore
it, an idea imported from masked modeling in language and vision. These methods differ in
what they mask. GraphMAE [22] masks node features and reconstructs them with a scaled cosine
error and a re-masking step, and shows that a plain autoencoder can then match or exceed
contrastive baselines. MaskGAE [29] masks structure instead, treating masked edge prediction
as the pretext and reconstructing the removed edges from the visible ones, with theory that
links the objective to contrastive learning. RARE [30] masks and predicts in the latent
space rather than the input space for added robustness. These objectives are attractive
because they need no negative sampling and carry over a recipe that has worked well
elsewhere.

Masking, though, is a pretext applied to data that is otherwise complete. When attributes are
genuinely missing, the feature-reconstruction signal largely overlaps with what neighborhood
aggregation already recovers, so it adds little to what the encoder computes on its own. Our
experiments bear this out. A feature-reconstruction objective plateaus early and gives no
downstream gain under missingness. Faithful reconstruction and a useful representation are
not the same goal in this regime.

### 2.3 Self-supervision for attribute-missing graphs

The methods closest to ours are SSL and multi-view models built for attribute-missing
graphs. AmGCL [14] works in two stages. A Dirichlet-energy-minimization precoder writes
information into the missing entries, and a self-supervised contrastive module then learns
representations by maximizing a lower bound on the mutual information of the latent code,
together with a structure-attribute energy-based reconstruction. AmGCL commits to a single
contrastive design, does not analyze which objective or view produces its gains, and is
evaluated only at small scale. MATE [7] takes a multi-view route, building complementary
views of the graph and enforcing agreement across them for imputation. Because supervision
never reaches its per-node augmented quantities directly, such designs are prone to
memorization, and shallow encoders can be starved of signal when missingness is high. MOBA
[8] shares the multi-view philosophy but sets out to fix these weaknesses, using a
collaborative strategy that cuts redundant information between views while keeping them
consistent. Its view design is still fixed in advance, it does not explain why a given view
or objective works, and it is not demonstrated at large scale. AIAE [9] is autoencoder-based
and reconstruction-centric, framed around imputation quality rather than downstream accuracy
under severe missingness. Outside the homogeneous setting, HGCA [11] joins attribute
completion and representation learning for heterogeneous graphs through an unsupervised
contrastive objective, completing missing attributes at a fine granularity with an augmented
network. Our own conference work [71] introduced a neighborhood-centric self-supervised
formulation, which the present study extends.

The pattern across these methods is consistent. Each fixes one objective and one view and
tunes it, and none investigates which self-supervised objective and which view help under
missingness, or why. Their results are also confined to small graphs and usually a single
missing rate, so behavior across the missingness range and at scale is left unexamined.

### 2.4 Adjacent communities

Two nearby literatures combine missingness with contrastive learning under a different task
or connectivity model. Incomplete multi-view clustering pairs absent information with
contrastive objectives but targets clustering rather than node classification. COMPLETER
[28] handles the views of an incomplete multi-view dataset with a contrastive-prediction
objective, and MCGC [52] learns a consensus graph with a contrastive regularizer because the
observed graph may be noisy or incomplete. On the higher-order side, SGHFP [51] propagates
features over hypergraphs to cope with missing attributes. We note these links but do not
treat them as central, since they differ from our problem in task or in graph model.

---

## 3. Positioning of Our Work

Seen together, the prior work leaves a clear opening. Imputation methods, whether statistical
propagation or learned generative models, optimize reconstruction instead of downstream
utility, and the learned variants do not scale. Imputation-free architectural methods and
per-node methods either add no self-supervisory signal or memorize node-specific quantities
that do not generalize. Structure-leaning methods exploit topology effectively but are built
for a broader joint-deficiency problem. General graph SSL is powerful but untested as a tool
for missing-feature learning, and its own literature shows that view construction is a
brittle, largely handcrafted axis. The SSL methods built for attribute-missing graphs each
settle on one objective and one view without asking which choices matter or why, and report
results only at small scale and a single missing rate.

Our work responds to this gap directly. We evaluate representations by downstream
node-classification accuracy under missingness, at a substantially larger scale than prior
attribute-missing studies and across a range of missing rates. The core of the paper is then
a systematic study of which self-supervised objectives and which views help under
missingness. The pattern we find has a clear explanation: stochastic, structure-based
objectives whose signal is not redundant with neighborhood aggregation are the effective
ones, while feature-reconstruction objectives are not. This effect is not specific to our
encoder. The objective transfers to other backbones, which indicates that the principle,
rather than any particular architecture, drives the gains.
